\clearpage
\chapter*{Erratum to Exposé no. 190}
\addcontentsline{toc}{chapter}{Erratum to Exposé no. 190}
\setcounter{footnote}{0}
\src{300}

\noindent\textbf{3. Descent technique and existence theorems in algebraic
geometry [TDTE I]. I: Generalities. Descent along faithfully flat morphisms.}

\begin{center}
[Bourbaki Seminar, vol. 12, 1959/60, no. 190, 29 pp.]
\end{center}

\noindent\textbf{Page 1, line 10.} --- It now seems excessive to say that the
descent technique is “the basis of most existence theorems in algebraic
geometry.” This is true to a large extent for the nonprojective techniques
treated in the first two exposés of the TDTE I--VI series, but not for the
projective techniques (TDTE IV, V, VI).

\noindent\textbf{Page 5, line 16.} --- It is unnecessary to assume that
$\alpha$ is an $\mathcal F$-descent morphism.

\noindent\textbf{Page 20, remark.} --- A quasi-finite, étale, surjective
morphism $S'\longrightarrow S$, or a morphism
$\operatorname{Spec}(\widehat A)\longrightarrow\operatorname{Spec}(A)$, need
not be a strict descent morphism, even if $A$ is the local ring of an algebraic
curve over an algebraically closed field $k$ and $S=\operatorname{Spec}(A)$.
Thus one can find a proper smooth morphism $f:X\longrightarrow S$ making $X$ a
principal bundle over $S$ under an elliptic curve $E$, such that
$f':X'\longrightarrow S'$ is projective while $f$ is not projective. This is
therefore also an example of a non-isotrivial homogeneous principal bundle
under an abelian scheme.

\noindent\textbf{Page 26, line 9.} --- Read “Chow--Lang” instead of
“Igusa--Lang.”

For various details concerning descent theory, see SGA VI, VII, VIII.

\par\smallskip
\begingroup\footnotesize
\endgroup
